\documentclass[11pt]{article}
% Кодировка и язык
\usepackage[utf8]{inputenc}           % Кодировка исходного файла (UTF-8)
\usepackage[T2A]{fontenc}             % Кодировка шрифтов для кириллицы
\usepackage[english,russian]{babel}   % Локализация: русские названия, переносы

% Поиск в PDF
\usepackage{cmap}                     

% Поля страницы
\usepackage{geometry}
\geometry{left=2cm, right=2cm, top=2cm, bottom=2cm}

% Математика
\usepackage{amsmath, amssymb, amsthm} % Базовые пакеты для математики
\usepackage{mathtools}                % Расширение amsmath (доп. символы и окружения)

% Картинки
\usepackage{graphicx}                 % Вставка картинок (\includegraphics)
\graphicspath{{pictures/}}            % Путь к папке с картинками
\usepackage{float}                    % Позволяет использовать [H] для жёсткого позиционирования
\usepackage{wrapfig}                  % Обтекание текста вокруг картинок

% Таблицы
\usepackage{array}                    % Расширенные возможности для таблиц
\usepackage{booktabs}                 % Красивые таблицы (\toprule, \midrule, \bottomrule)
\usepackage{multirow}                 % Объединение строк в таблицах
\usepackage{tabularx}                 % Таблицы с автоматической шириной колонок
\usepackage{ragged2e}                 % Улучшенное выравнивание (используется в tabularx)

% Списки
\usepackage{enumitem}                 % Гибкая настройка списков (itemize, enumerate)

% Колонтитулы
\usepackage{fancyhdr}                 

% Гиперссылки
\usepackage{xcolor}                   % Работа с цветами
\usepackage{hyperref}                 % Гиперссылки, оглавление, перекрёстные ссылки

% Умные ссылки
\usepackage{cleveref}                 

% Типографика
\usepackage[expansion=false]{microtype}                % Улучшает переносы и межсловные расстояния

% Подписи к картинкам и таблицам
\usepackage{caption}                  % Настройка подписей (шрифт, отступы)

% Рисование схем и графиков
\usepackage{tikz}                     % Рисование схем, графиков, диаграмм
\usepackage{pgfplots}                 % Построение графиков функций (на основе tikz)
\pgfplotsset{compat=newest}           % Версия совместимости pgfplots

% Алгоритмы и код
\usepackage{algorithm}                % Окружения для алгоритмов
\usepackage{algorithmic}              % Синтаксис для алгоритмов
\usepackage{listings}                 % Вставка кода с подсветкой синтаксиса


% Настройка нумерации секций и подсекций
\renewcommand{\thesection}{\arabic{section}}                   
\renewcommand{\thesubsection}{\thesection.\arabic{subsection}} 


% Настройка списков: убираем лишние отступы слева
\setlist[itemize]{leftmargin=*}
\setlist[enumerate]{leftmargin=*}


% Определение окружений для теорем, лемм, определений и т.д.
% Нумерация привязана к номеру секции
\newtheorem{theorem}{Теорема}[section]
\newtheorem{definition}{Определение}[section]
\newtheorem{remark}{Замечание}[section]
\newtheorem{lemma}{Лемма}[section]
\newtheorem{example}{Пример}[section]
\newtheorem{consequence}{Следствие}[section]
\newtheorem{algorithmm}{Алгоритм}
\newtheorem{question}{Вопрос}[section]
\newtheorem{answer}{Ответ}[section]
\newtheorem{statement}{Утверждение}[section]


% Определение цветов
\definecolor{pblue}{rgb}{0.13,0.13,1}    
\definecolor{pgreen}{rgb}{0,0.5,0}       
\definecolor{pred}{rgb}{0.9,0,0}         
\definecolor{pgrey}{rgb}{0.46,0.45,0.48}

\DeclareMathOperator{\SAT}{SAT}
\DeclareMathOperator{\tSAT}{3-SAT}
\DeclareMathOperator{\tCOL}{2COL}
\DeclareMathOperator{\trCOL}{3COL}
\DeclareMathOperator{\VERTEXCOVER}{VERTEXCOVER}
\DeclareMathOperator{\DHAMPATH}{DHAMPATH}
\DeclareMathOperator{\CLIQUE}{CLIQUE}
\DeclareMathOperator{\INDSET}{INDSET}
\DeclareMathOperator{\Poly}{P}
\DeclareMathOperator{\NotDet}{N}
\DeclareMathOperator{\NP}{NP}
\DeclareMathOperator{\coNP}{coNP}
\DeclareMathOperator{\PH}{PH}
\DeclareMathOperator{\poly}{poly}
\DeclareMathOperator{\EXACTSETC}{EXACTSETCOVER}
\DeclareMathOperator{\SUBSUM}{SUBSET-SUM}
\DeclareMathOperator{\KNAPSACK}{KNAPSACK}

\newcommand{\lectureNumber}{6}
\newcommand{\lectureTitle}{$\NP$-полнота задачи $\EXACTSETC$. Не $\NP$-полные вариации задачи о рюкзаке. Теорема Махэни. Полиномиальная иерархия}
\newcommand{\courseName}{Формальные языки и сложность вычислений}
\newcommand{\semester}{3 семестр, 2026/2027 уч. год}

\pagestyle{fancy}
\fancyhf{}
\fancyhead[L]{Я слишком ленивый, чтобы держать на тебя злость}
\fancyhead[R]{\thepage}
\setlength{\headheight}{15pt}

\hypersetup{
    colorlinks=true,
    linkcolor=blue!70!black,
    citecolor=blue!70!black,
    filecolor=blue!70!black,
    urlcolor=blue!70!black,
}

\begin{document}


\begin{titlepage}
    \centering
    
    {\large Московский физико-технический институт \par}
    {\large Высшая школа программной инженерии \par}
    \vspace{2cm}

    {\Large \textbf{\courseName} \par}
    {\large \semester \par}
    
    \vspace*{\fill}

    {\Huge \textbf{Лекция \lectureNumber} \par}
    {\Huge \textbf{\lectureTitle} \par}
    
    \vspace*{\fill}
    
    {\large Лектор: Зухба Анастасия Викторовна <3\par}
    {\large Автор: \href{https://t.me/sol_m_07}{Соля} \par}

\end{titlepage}

\pagenumbering{arabic}

\section{$\NP$-полнота задачи $\EXACTSETC$}

Рассмотрим следующие три $\NP$-полные задачки:

\[\EXACTSETC = \{(n, S_1, \ldots, S_m) \; | \; \text{можно выбрать непересекающиеся множества}\]
\[S_{i_1}, \ldots, S_{i_k}, \text{ в объединении дающие все числа от 1 до } n\}\]
\[\SUBSUM = \{(k_1, k_2, \ldots, k_m, N) \; | \; \exists a \in \{0, 1\}^n \ \sum_{i = 1}^m a_i k_i = N\}\]
\[\KNAPSACK = \{((w_1, \ldots, w_n), (v_1, \ldots, v_n), W, V) \; | \; x_i \in \{0, 1\}, \ \sum_{i = 1}^n x_i w_i \leq W, \ \sum_{i = 1}^n x_i v_i \geq V\}\]

\begin{remark}
  Задача о рюкзаке имеет множество различных постановок. Многие из них (но не все) 
  являются $\NP$-полными.
\end{remark}

\begin{proof}
  Докажем $\NP$-полноту задачи $\EXACTSETC$ с помощью полиномиальной сводимости из задачи $\tSAT$ (принадлежность $\EXACTSETC$ к $\NP$ очевидна, 
  так как просто берём сертификат и жизнь сразу налаживается). 
  Пусть дана формула 
  \[\varphi = \underbrace{(x \vee \overline{y} \vee z)}_{A} \wedge 
  \underbrace{(\overline{x} \vee z \vee \overline{t})}_{B} \wedge 
  \underbrace{(\ldots)}_{C} \wedge {\ldots}\]
  в 3-КНФ. Мы хотим свести её к задаче $\EXACTSETC$. Дадим каждому дизъюнкту ``имя'' ($A, B, C, \ldots$). 
  Тогда множество, которое мы хотим получить в задаче $\EXACTSETC$, будет включать в себя
  \begin{itemize}
    \item $x, x_0, x_1$ для каждой переменной $x$
    \item $C, C_1, C_2, C_3$ для каждого дизъюнкта $C$
  \end{itemize}
  Из каких множеств мы будем его составлять? Сначала введём ``заглушки'':
  \begin{itemize}
    \item $\{x, x_0\}$, $\{x, x_1\}$ для каждой переменной $x$
    \item $\{C\}$, $\{C, C_1\}$, $\{C, C_2\}$, $\{C, C_3\}$, $\{C, C_1, C_2\}$, $\{C, C_1, C_3\}$, $\{C, C_2, C_3\}$ для каждого дизъюнкта $C$
  \end{itemize}
  Далее введём множества, которые могут сначала показаться контринтуитивными (захочется поменять местами 0 и 1), но на самом деле всё окей:
  \begin{itemize}
    \item $\{x_0, C_j, \ldots\}$ для каждой переменной $x$, если $x$ встречается в дизъюнкте $C$ в позиции $j$ \textbf{без отрицания}
    \item $\{x_1, D_k, \ldots\}$ для каждой переменной $x$, если $x$ встречается в дизъюнкте $D$ в позиции $k$ \textbf{с отрицанием}
  \end{itemize}
  Если $\varphi$ выполнима, то возьмём выполняющий набор. Для каждой переменной $x$ будем брать 
  \begin{itemize}
    \item $\{x_0, C_j, \ldots\}$ и ``заглушку'' $\{x, x_1\}$, если $x$ в выполняющем наборе равна 1
    \item $\{x_1, D_k, \ldots\}$ и ``заглушку'' $\{x, x_0\}$, если $x$ в выполняющем наборе равна 0
  \end{itemize}
  После этого для каждого дизъюнкта $C$ возьмём подходящую ``заглушку'', в которой будет $C$ и все не вошедшие в объединение $C_j$. 
  Тут важно заметить, что заглушка $\{C, C_1, C_2, C_3\}$ действительно не понадобится и всё пройдёт корректно, так как в противном случае делаем вывод о том, 
  что ни одна переменная не ``заединичила'' дизъюнкт $C$, что противоречит тому, что мы взяли выполняющий набор.
  Посмотрим же, почему тогда мы решили задачу $\EXACTSETC$:
  \begin{enumerate}
    \item По построению мы охватили все элементы $x, x_0, x_1$
    \item По построению мы охватили все элементы $C, C_1, C_2, C_3$
    \item Ни один элемент переменной не был учтён дважды (также очевидно из построения)
    \item Ни один элемент дизъюнкта не был учтён дважды, о чём давайте далее чуть подробнее
  \end{enumerate}
  Почему последнее верно? Давайте просто заметим, что $C_j$ по построению возникнет только тогда, когда $x$ встретится в дизъюнкте $C$ в позиции $j$ (с отрицанием или без). 
  Мы для каждой переменной $x$ берём только один набор вида $\{x_\alpha, C_j, \ldots\}$. Если мы получили $C_j$ дважды, то нашлись наборы $\{x_\alpha, C_j, \ldots\}$ и $\{y_\beta, C_j, \ldots\}$.
  Но ведь в дизъюнкте $C$ не может быть двух разных переменных в одной позиции, значит $x = y$. Но тогда мы взяли два набора для одной переменной, что невозможно.
  Ура, уже доказано
  \[\varphi \in \tSAT \implies f(\varphi) \in \EXACTSETC\]
  В обратную сторону всё даже легче. У нас есть какая-то 3-КНФ $\varphi$ и построенный по ней $f(\varphi)$ --- набор множеств, который даёт решение задачи $\EXACTSETC$.
  Рассмотрим тогда это точное покрытие. Для каждой переменной $x$ оно содержит либо $\{x, x_0\}$, либо $\{x, x_1\}$, так как $x$ есть только там. Тогда оставшаяся $x_\alpha$ покрывается набором 
  вида $\{x_\alpha, C_j, \ldots\}$. На самом деле уже на этом шаге мы для каждого дизъюнкта $C$ покрыли один из его элементов $C_j$. Ведь иначе нам бы пришлось покрыть $\{C, C_1, C_2, C_3\}$, 
  что невозможно сделать ``заглушками'' (одной большой нет, а между собой они пересекаются как минимум по $C$). Тогда просто посмотрим на попавшие в покрытие наборы вида $\{x, x_\alpha\}$ и 
  положим $x = \alpha$ для каждой переменной $x$. Из рассуждений выше как раз становится ясно, что полученный набор значений переменных будет выполняющим для $\varphi$. Ура, доказали и обратное:
  \[f(\varphi) \in \EXACTSETC \implies \varphi \in \tSAT\]
  Осталось только пояснить, почему $f$ --- полиномиально вычислимая по длине $\varphi$. Действительно, количество переменных --- полином от длины $\varphi$, количество дизъюнктов --- полином от длины $\varphi$. 
  Всё строится за полином от этих количеств, значит за полином от длины $\varphi$. Ура, доказали $\NP$-полноту задачи $\EXACTSETC$.
\end{proof}

\section{$\NP$-полнота задачи $\SUBSUM$}

\begin{proof}
  Докажем $\NP$-полноту задачи $\SUBSUM$ с помощью полиномиальной сводимости из задачи $\EXACTSETC$. Пусть у нас есть экземпляр задачи $\EXACTSETC$ --- $\{n, S_1, S_2, \ldots, S_m\}$.
  Тогда в $(m+1)$-ричной системе счисления построим числа $k_1, k_2, \ldots, k_m$ и $N$ следующим образом:
  \begin{itemize}
    \item $k_i$ --- число, у которого в $j$-той цифре стоит 1, если $j \in S_i$, и 0 иначе
    \item $N$ --- число из $n$ единиц
  \end{itemize} 
  Тогда очевидно, что если мы можем выбрать непересекающиеся множества $S_{i_1}, \ldots, S_{i_h}$, то сумма соответствующих чисел $k_{i_1} + \ldots + k_{i_h}$ будет равна $N$.
  В обратную сторону, если мы можем выбрать числа $k_{i_1}, \ldots, k_{i_h}$, сумма которых равна $N$, то в каждой цифре суммы будет стоять 1, что означает, что выбранные множества $S_{i_1}, \ldots, S_{i_h}$ непересекающиеся и дают в объединении все числа от 1 до $n$.
  Ура, доказали $\NP$-полноту задачи $\SUBSUM$ (полиномиальная вычислимость функции $f$ очевидна, так как мы просто строим числа в $(m+1)$-ричной системе счисления).
\end{proof}

\section{Задача о рюкзаке}

\begin{proof}
  Докажем $\NP$-полноту задачи $\KNAPSACK$ с помощью полиномиальной сводимости из задачи $\SUBSUM$. Пусть у нас есть экземпляр задачи $\SUBSUM$ --- $\{k_1, k_2, \ldots, k_m, N\}$.
  Тогда просто положим $w_i = v_i = k_i$, $W = V = N$. Тогда очевидно, что если мы можем выбрать числа $k_{i_1}, \ldots, k_{i_h}$, сумма которых равна $N$, то мы можем положить в рюкзак предметы с весами и ценностями $w_{i_1}, \ldots, w_{i_h}$ и получить суммарный вес и ценность равные $N$.
  В обратную сторону, если мы можем положить в рюкзак предметы с весами и ценностями $w_{i_1}, \ldots, w_{i_h}$ и получить суммарный вес и ценность равные $N$, то мы можем выбрать числа $k_{i_1}, \ldots, k_{i_h}$, сумма которых равна $N$.
  Ура, доказали $\NP$-полноту задачи $\KNAPSACK$ (полиномиальная вычислимость функции $f$ очевидна, так как мы просто копируем числа).
\end{proof}

\subsection{Задача о непрерывном рюкзаке}

На самом деле не все вариации задачи о рюкзаке являются $\NP$-полными. Например, если разрешить дробные значения $x_i$, то задача становится полиномиально решаемой. 
Действительно, в этом случае можно просто отсортировать предметы по убыванию ценности на единицу веса и брать их (опционально дробно) в рюкзак до тех пор, пока не достигнем ограничения по весу. 
Если набрали нужную ценность, то задача решена.

\subsection{А что если $W$ записана в унарной системе счисления?}

Если $W$ записана в унарной системе счисления, то в этом случае задача остается полиномиально решаемой, поскольку можно использовать динамическое программирование:
\begin{algorithm}
\caption{Задача о рюкзаке}
\begin{algorithmic}[1]
    \STATE \(dp[0..W] \gets 0\)
    \FOR{все предметы (w, v)}
        \FOR{\(j = W\) \textbf{ downto } \(w\)}
            \STATE \(dp[j] \gets \max(dp[j], dp[j - w] + v)\)
        \ENDFOR
    \ENDFOR
    \STATE \textbf{return} \(dp[W]\)
\end{algorithmic}
\end{algorithm} \leavevmode
\\
Получили сложность $O(nW)$ и доп память $O(W)$.

\section{Теорема Махэни}

\begin{definition}
  \textbf{Разреженный язык} --- это формальный язык, для которого количество слов длины $n$ ограничено полиномом от $n$.
\end{definition}

\begin{theorem}[Махэни]
  Если существует язык, который одновременно является $\NP$-полным и разреженным, то $\Poly = \NP$.
\end{theorem}

\begin{lemma}
  Любой унарный язык является разреженным.
\end{lemma}

\begin{consequence}
  Если существует унарный $\NP$-полный язык, то $\Poly = \NP$.
\end{consequence}

\begin{remark}
  Следствие выше как раз и позволяет нам понять, что перекодировать $\NP$-полный язык в унарную запись --- не супер идея.
\end{remark}

\begin{proof}
  Пусть $S$ --- $\NP$-полный и разреженный язык. Так как он $\NP$-полный, то $\SAT$ полиномиально сводится к нему, т.е. 
  \[\exists f \text{ --- всюду определенная полиномиально вычислимая функция: } x \in \SAT \iff f(x) \in S\]
  Покажем тогда, как решить $\SAT$ за полиномиальное время: \\ \\
  У нас точно будет монотонный полином $p(m)$, ограничивающий количество слов длины не более $m$ 
  (изначально в определении сказано ровно про $m$ и нет ничего про монотонность, но первое очевидно следует из того, 
  что есть полином, ограничивающий количество слов длины ровно $m$, и что мы можем просто взять его максимум на отрезке $[0, m]$ 
  и домножить на $m$, а второе следует из того, что любой не монотонный полином можно сверху ограничить монотонным). \\ \\
  Теперь рассмотрим формулу $\varphi$ длины $N$ и в ней $n$ переменных. При подстановке некоторых значений переменных мы 
  получаем формулу $\widetilde{\varphi}$ длины не более $c \cdot (N + 1)$, где $c$ --- константа. Подставляя по одному значению переменной, 
  получаем каждый раз увеличение кол-ва формул в $2$ раза. Мы знаем, что выполнимость $\varphi$ на самом деле равносильна выполнимости 
  хотя бы одной из формул $\varTheta_{i}$, полученных после подстановки. Пока что нас конечно не радует, что их $2^n$, поэтому посмотрим, 
  как их отсекать без потери информации. \\ \\
  У нас сейчас есть список формул $\varTheta_{1}, \ldots, \varTheta_{k}$, каждая из которых имеет длину не более $c \cdot (N + 1)$. 
  Если взять дизъюнкцию двух из них $\varTheta_{i} \vee \varTheta_{j}$, то она будет иметь длину не более $c' \cdot (N + 1)$. 
  Теперь вспомним, что $f$ --- полиномиально вычислимая функция $\implies f(\varTheta_{i} \vee \varTheta_{j})$ будет ограничена 
  сверху полиномом по $N$. Пусть этот полином --- $R(N)$. Тогда введём $T(N) = p(R(N))$. Получается, что среди $S$ не более $T(N)$ 
  слов длины не более $R(N)$. \\ \\
  Пока $k > T(N) + 1$ и мы ещё не подставили все значения переменных, подставляем очередное значение переменной, 
  после чего берём любую формулу (например $\varTheta_{1}$) и дизъюнкцию её со всеми другими формулами. Положим 
  \[z_i = f(\varTheta_{1} \vee \varTheta_{i})\]
  \paragraph{Случай 1}
  \[\exists i \neq j: \ z_i = z_j\]
  Рассмотрим тогда формулы $\varTheta_{1} \vee \varTheta_{i}$ и $\varTheta_{1} \vee \varTheta_{j}$. Так как $f$ --- всюду определённая функция,
  то $z_i = z_j \implies f(\varTheta_{1} \vee \varTheta_{i}) \in S \iff f(\varTheta_{1} \vee \varTheta_{j}) \in S$.
  Тогда мы можем просто удалить из списка формулу $\varTheta_{j}$, так как она не несёт никакой новой информации.
  \paragraph{Случай 2}
  \[\forall i \neq j: \ z_i \neq z_j\]
  Тогда из разреженности языка $S$ следует, что $\exists i: z_i \notin S \implies \varTheta_{1}$ не выполнима и её можно удалить из списка. \\ \\
  Так мы каждый раз удаляем, пока не вернём $k \leq T(N) + 1$. Затем просто подставляем очередное значение переменной (если ещё не завершили), 
  получая $k \leq 2T(N) + 2$. \\ \\
  Оказывается, что всё это мы сделали за полиномиальное время по $N$, но ведь на самом деле мы решили $\SAT$. В полученном списке формул будут лишь 
  истина и ложь, так как мы подставили все значения переменных. Тогда просто проверяем, есть ли среди них истина.
\end{proof}

\begin{remark}
  Унарный $\NP$-трудный язык можно построить, например, с помощью функции трудолюбия Радо.
\end{remark}

\section{Полиномиальная иерархия}

Вспомним пройденное и адаптируем под новые реалии:

\[\Sigma_{k}^{p}: \ x \in A \iff \exists y_1 \forall y_2 \ldots Q_k y_{k} \ V(x, y_1, \ldots, y_k) = 1\]
\[\Pi_{k}^{p}: \ x \in A \iff \forall y_1 \exists y_2 \ldots Q_k y_{k} \ V(x, y_1, \ldots, y_k) = 1\]
$V$ вычислим за полиномиальное время; $\left\lvert y_i \right\rvert \leq p(\left\lvert x \right\rvert)$, где $p$ --- полином.
\[\Sigma_{0}^{p} = \Pi_{0}^{p} = \Poly; \quad \Sigma_{1}^{p} = \NP; \quad \Pi_{1}^{p} = \coNP\] 
\[\Sigma_{k}^{p} \cup \Pi_{k}^{p} \subseteq \Sigma_{k+1}^{p} \cap \Pi_{k+1}^{p}; \quad \PH = \bigcup_{k \geq 0} \Sigma_{k}^{p} = \bigcup_{k \geq 0} \Pi_{k}^{p}\]

\section{Дисклеймер для преподавателей перед любой сложной темой}

\begin{figure}[h]
  \centering
  \includegraphics[width=0.8\linewidth]{Опоссумы.jpg}
\end{figure}

\end{document}
